\section{\label{apx:extended}Extended design and experimental setup}
This appendix carries the formal development and implementation detail compressed out of \S\ref{sec:design} and \S\ref{sec:setup}.

\subsection{\label{apx:extdesign}Design}
Let $\theta_0$ be a small instruction-tuned open-weight LLM.
Each run applies a privacy mechanism $m_i$ at a given strength,
acting on either the adaptation data or the adaptation procedure,
and produces an adapted model.
Each run yields two measurements.
\textbf{Semantic distortion} $S_i$ is the mean cosine distance between sentence embeddings of the original and the rewritten training text;
where the mechanism leaves the text intact,
it is measured between generations of the adapted and the base model on a fixed probe set.
$S$ is computed without reference to the privacy parameter that produced it.
\textbf{Alignment degradation} $Y_i$ is the drop from the same model's non-private baseline on the evaluation suite of \S\ref{sec:eval}.

The two candidate pathways make opposite predictions about what survives conditioning on distortion:
if degradation runs through meaning,
the association between privacy level and $Y$ vanishes once $S$ is held fixed;
if it does not,
a level effect remains at fixed $S$. Splitting the association that way and writing it linearly yields
\begin{equation}
  Y_i = a_{m_i} + b_{m_i} d_i + c S_i + e_i,
  \label{eq:modelext}
\end{equation}
where $d_i$ is the privacy-level rank of run $i$ within its own mechanism,
$a_m$ and $b_m$ are a per-mechanism intercept and level slope,
respectively,
and $e_i$ is the error term.
A random intercept per base model is added to $a_m$.
$Y$,
$d$ and $S$ are z-scored,
so $b_m$ and $c$ are standardized effects on a common scale.
Baseline runs have $Y \equiv 0$ by definition:
they set the zero point and are not regression rows.

Despite the mechanism index on $a_m$ and $b_m$,
\eqref{eq:model} is one regression estimated jointly over all mechanism runs through indicator variables:
seven coefficients fitted in a single pass,
with the per-model random intercept as a mixed-effects term.
The pre-registered cluster-robust correction by model~$\times$~mechanism cell is recorded as a deviation in \apx{deviations}.
Each row activates the intercept and slope of its own mechanism only.
The $S$ column carries no indicator,
so $c$ is shared and pools mechanisms that disagree about how much distortion a given privacy level buys.
Joint estimation puts $\hat c$ and each $\hat b_m$ in one covariance matrix,
so comparing them is a Wald test on linear combinations of coefficients.

The privacy level enters \eqref{eq:model} only through its rank $d_i$.
Within each mechanism the settings are ordered from strictest to loosest,
coded $1$ upward,
and z-scored,
so mechanisms with different numbers of levels enter on the same scale.
The nominal budgets are incomparable across families:
private optimization gives example-level $(\varepsilon, \delta)$-DP~\citep{abadi2016deep},
sanitization word-level metric local differential privacy (LDP)~\citep{chatzikokolakis2013broadening}.
Rank coding rests on the ordering of the settings alone,
which leaves $S$ as the one regressor common to every mechanism.

We state the claim in two stages,
because a shared coefficient can be statistically distinct from zero and still too small to matter.
The first stage fixes a negligibility floor $\Delta = 0.10$ in standardized units,
committed before the main pass together with the primary $S$ metric and the exclusion rules,
and asks whether $|c| \ge \Delta$:
distortion moves alignment by an amount worth acting on.
The value of $\Delta$ is arbitrary and the conclusion does not turn on it;
it marks the point below which an effect is negligible,
so a confidence interval for $c$ contained in $(-\Delta, \Delta)$ settles the first stage the other way,
establishing the practical absence of the semantic pathway by TOST.
The second stage is the substantive one:
$|c| \gg |b_m|$ for every mechanism,
distortion predicting degradation far more strongly than the nominal privacy level,
tested as $|c| > |b_m|$ on the joint covariance of \eqref{eq:model}.

Within a single mechanism $S$ and $\varepsilon$ are correlated by construction,
so no one mechanism separates $b_m$ from $c$. The design produces variation in $S$ at fixed privacy from two sources.
The first is the contrast between families,
which reach comparable privacy with different amounts of semantic damage:
sanitization passes out-of-vocabulary tokens through untouched while random replacement destroys them,
and personally identifiable information (PII) is typically out of vocabulary (\S\ref{sec:mech}).
The second is a non-private control.
If it degrades alignment as much as the private conditions at equal $S$,
semantic damage rather than privacy is the cause;
if it does not,
the residual measures mechanism-specific,
non-semantic effects.
We report VIF for $S$,
falling back to the within-level cross-mechanism contrast,
which is free of $\varepsilon$ -- $S$ confounding.

\subsection{\label{apx:extsetup}Experimental setup}
Here, we describe datasets,
privacy mechanisms,
models and evaluation suite.

\subsection{Datasets}
The text the mechanisms perturb stands in for the sensitive data.
The two datasets differ in \textbf{entity density},
because privacy risk concentrates in names,
numbers and facts,
and the mechanisms treat these tokens differently from common words:
rare tokens occupy sparse regions of embedding space,
so sanitization replaces them with distant,
fact-changing substitutes,
and private optimization suppresses the memorization of exactly such rare sequences.

\paragraph{Entity-rich.}
\texttt{nvidia/Nemotron-PII}~\citep{nvidia2025nemotronpii} is a synthetic dataset of emails,
clinical notes,
forms and free text,
with gold span-level PII annotations.
The gold spans make the entity component of $S$ exact rather than dependent on a noisy named entity recognition model,
and let us measure whether an adapted model reproduces or corrupts specific entities.
We restrict to the biomedical domains for simplicity.
The mechanisms perturb text without regard to domain,
so one domain suffices to study them and a broader set would add cost without changing what we measure.
We keep records with at least one parsable span.

Each document is rendered once into an instruction example,
reused by every condition,
under one of three templates sampled per record:
extract a labeled entity,
draft a document of the given type and domain,
or classify the document type.
The training set is $16{,}384$ rows.
A further $375$ records are held out for the entity-fidelity probe,
disjoint \textit{by source document} rather than by rendering,
since one document appears under several renderings and splitting on rows would leak.
Another $375$ are sampled from the training rows for the leakage probe;
these are not held out,
because the probe asks whether training memorized them.

\paragraph{Entity-poor.}
A fixed subset of Alpaca~\citep{taori2023alpaca} of the same size:
generic assistant tasks with few named entities or numbers,
and the standard small instruction-tuning dataset,
which keeps results comparable to prior work.
Gold spans define the entity and carrier components of $S$,
so those two components are reported on the entity-rich dataset,
and the entity-poor dataset carries the composite $S$ alone.
Dataset construction and perturbation run as a single pass emitting both datasets together,
so the perturbed datasets are matched on both sides of the entity-density contrast.

\subsection{\label{apx:mech}Privacy mechanisms}
The registered grid has four conditions:
a non-private baseline,
a non-private distorting control,
and two mechanisms carrying formal guarantees.

\paragraph{No privacy (baseline).}
Standard LoRA fine-tuning on the clean dataset,
defining the zero point of both $Y$ and $S$.

\paragraph{Random token replacement.}
Text is tokenized on whitespace,
numbers and identifiers included,
and each token is independently replaced with probability $p \in \{0.05, 0.15, 0.30, 0.50\}$ by a token drawn uniformly from its own dataset vocabulary.
A higher rate is the stricter setting,
so the rank order is reversed relative to the $\varepsilon$-indexed mechanisms.
It carries no privacy guarantee;
its purpose is to vary distortion -- entity distortion in particular,
since whitespace tokenization exposes identifiers -- independently of any privacy budget.

\paragraph{Sanitization.}
We use TEM~\citep{carvalho2021tem} rather than a plain SanText-style exponential mechanism over the full vocabulary~\citep{yue2021differential}.
TEM restricts the candidate set to words within a distance threshold and folds the remaining mass into a truncation symbol,
avoiding the plain mechanism's low-probability substitutions drawn from arbitrarily distant words.
For a word $x$ with embedding $e_x$,
the near set is $N = \{w : \lVert e_x - e_w \rVert_2 \le \gamma\}$ and the far set $F = V \setminus N$,
with utility $u(w) = -\lVert e_x - e_w \rVert_2$ and a truncation symbol $\bot$ of utility $u(\bot) = -\gamma + \tfrac{2}{\varepsilon} \ln|F|$;
a substitute is drawn from $N \cup \{\bot\}$ with probability proportional to $\exp(\tfrac{\varepsilon}{2} u)$,
and if $\bot$ is drawn a word is emitted uniformly from $F$. We use word-level $\varepsilon \in \{0.5, 1, 2, 3, 4, 5, 6, 12\}$ and $\gamma = 5.0$,
sampling per token occurrence but memoizing the per-word-type distribution across occurrences and datasets,
which makes the mechanism tractable over a $400$k-word vocabulary.

Word distances come from pre-trained GloVe 6B.300d vectors~\citep{pennington2014glove}.
Perturbation runs once and its output is consumed by all three models;
deriving the metric from a base model's own embedding matrix would make the sanitized dataset model-specific and entangle $S$,
the pooled regressor on which the design rests,
with model identity.
GloVe is also a \textit{word}-level vocabulary,
matching the units in which the entity component of $S$ is defined,
where a mechanism over BPE tokens would fragment entities across sub-tokens.
Out-of-vocabulary passthrough follows from the same choice:
medical record numbers,
phone numbers and email addresses are out of vocabulary in GloVe,
so TEM leaves them intact and \textit{preserves} literal entities where random replacement destroys them at matched nominal privacy.
That asymmetry is the identification strategy of \S\ref{sec:design},
and subword vectors such as fastText would remove it.

Both data-side mechanisms rewrite only the content fields of each instructionized example,
leaving the instruction template intact:
the instruction is a fixed task specification rather than sensitive data,
and perturbing it corrupts what the model is asked to do instead of the signal the mechanism acts on.
Perturbed rows keep the original schema and gold-span offsets,
so entity-level distortion is computed by joining each perturbed row to its clean counterpart on the example identifier.

\paragraph{Private optimization.}
DP-SGD~\citep{abadi2016deep} is applied to the LoRA parameters at example-level $\varepsilon \in \{1, 4, 8, 16\}$ with $\delta = 10^{-5}$.
Lots are drawn by Poisson subsampling,
each example included independently with probability $q = B/N$.
This is not interchangeable with shuffled fixed-size minibatches:
the amplification result the accountant applies is proved for Poisson sampling~\citep{mironov2019sgm},
and under shuffled batches the reported $\varepsilon$ would misstate the mechanism that was run.
Per-example gradients are clipped to $\lVert g \rVert_2 \le C$ and accumulated in fp32,
Gaussian noise $\mathcal{N}(0, (\sigma C)^2)$ is added to the sum,
and the result is divided by the \textit{configured} lot size rather than the number of examples drawn,
as Poisson sampling requires.
The private loop is written directly against the adapter parameters,
but privacy accounting is delegated to an established implementation:
a numerical inversion for the $\varepsilon \mapsto \sigma$ calibration,
and a R\'enyi accountant~\citep{mironov2017renyi} for the realized $\varepsilon$,
recomputed at the chosen $\sigma$ over the actual step count and logged as the reported guarantee.
An error in a training loop shows up in the loss;
an error in a hand-vendored accountant silently under-reports $\varepsilon$.
\apx{dpsgd} gives the implementation details,
the end-to-end validation of the calibration,
and the pitfalls specific to this branch.
In the registered grid the mechanisms are kept separate:
the private branch trains only on clean text.
A post-hoc block lifts that restriction (\S\ref{sec:scope}).
A post-hoc block lifts that guard to measure their interaction (\S\ref{sec:scope}).

\paragraph{Entity-targeted sanitization (post-hoc control).}
This variant was added after the main fit to separate the entity and carrier components of $S$,
which uniform sanitization moves together (\S\ref{sec:carrier}).
It applies TEM to the gold entity spans alone -- resampling each distinct entity once and substituting it at every occurrence -- and leaves the carrier text verbatim.
It runs at two budgets on the entity-rich dataset only,
since the targeting requires the gold spans.
At matched entity displacement it inflicts a fraction of the carrier damage that uniform sanitization does,
isolating the two components the uniform mechanism confounds.

\subsection{Models and adaptation}
We use three dense open-weight instruction-tuned models in the $1$ -- $2$B range:
Qwen3-1.7B~\citep{yang2025qwen3},
Llama-3.2-1B-Instruct~\citep{grattafiori2024llama},
and Falcon3-1B-Instruct~\citep{falcon3}.
The size is set by the grid:
with $324$ adapters to train and evaluate,
each must cost tens of GPU-minutes for the design to remain affordable.
Three \textit{families} rather than three checkpoints of one family let the per-model random intercept absorb base-model idiosyncrasy;
the claims of \S\ref{sec:design} require variation across pretraining recipes and tokenizers.
All three expose the same query and value projections,
keeping the adapter target set identical across conditions.
We exclude the newest hybrid small models -- gated linear attention with sparse mixture-of-experts -- since per-example gradient computation is awkward there and an accounting error on a novel architecture would invalidate the entire private column.

Every condition uses the same adaptation procedure:
LoRA~\citep{hu2022lora} with rank $16$,
$\alpha = 32$,
adapters on the query and value projections,
no dropout and no bias terms,
the base model loaded in bf16 and frozen,
and loss taken on completion tokens only using the same chat-template call the evaluation harness uses.
This gives $3{,}211{,}264$ trainable parameters.
Sequences are truncated at $2048$ tokens and training runs for $3$ epochs at a lot size of $256$,
with a warmup--stable--decay learning-rate schedule~\citep{hu2024minicpm} specified in \textit{fractions} of total steps,
so its shape is identical in every cell.
The conditions differ only in which dataset they consume,
except the private one,
which additionally runs the private loop,
gated on the presence of a target $\varepsilon$ rather than on a global flag.

Holding the procedure fixed leaves the mechanism as the only varying factor,
a precondition for \eqref{eq:model} to attribute differences in $Y$ to $S$ and the privacy level rather than to training hyperparameters.
The fixed learning rate carries a cost,
since private conditions often prefer a higher rate than clean ones and the shared value was tuned on the clean condition.
We measured it:
the tuned rate improves the private mechanism as well,
at every budget tested.
\apx{calib} reports the calibration runs,
the seed-noise floor that terminates them,
and two settings whose response is flat.

Epochs are held at $3$ because they affect privacy:
$\sigma$ is calibrated over the total step count,
so more epochs mean more composition and a noisier private mechanism at the same $\varepsilon$.
Training longer would give the non-private conditions more optimization and the private mechanism both more optimization and more noise,
an asymmetry that would enter $\hat b_m$ and confound the effect the regression measures.
Lot size is free of this coupling:
the private branch performs one backward pass per \textit{example},
so total work is independent of $B$,
and lot size is the parameter used to raise the gradient signal-to-noise ratio (\apx{dpsgd}).

\subsection{\label{apx:eval}Evaluation}
Every component of $Y$ is judge-free:
no LLM judge enters the measurement chain,
so the dependent variable is independent of any judge model's own alignment properties.
Decoding is greedy and capped at a fixed number of new tokens,
the same cap in every condition.

\paragraph{Truthfulness.}
TruthfulQA~\citep{lin2022truthfulqa} MC1 and MC2 on the validation split.
We score each option by the summed token log-probabilities of the option conditioned on the question prompt,
giving a likelihood over the fixed options.

\paragraph{Instruction following.}
IFEval~\citep{zhou2023instruction},
whose prompts carry programmatically verifiable constraints;
we report prompt-level and instruction-level accuracy using the reference verifiers.

\paragraph{Entity fidelity (hallucination).}
The model is asked to extract a gold-annotated entity from each of the $375$ held-out records,
with the requested label drawn deterministically per record.
Responses are scored against the gold spans into three exclusive outcomes:
\texttt{correct},
\texttt{corrupted} (every returned value occurs somewhere in the record, but the set is incorrect),
and \texttt{fabricated} (some returned value does not occur in the record at all, or nothing is returned).
The distinction separates the two pathways:
corruption indicates a model trained on distorted entities,
fabrication one that lost grounding entirely.

\paragraph{Refusal.}
Refusal is measured on the $520$ harmful instructions of AdvBench~\citep{zou2023universal} and on XSTest~\citep{rottger2024xstest},
whose $450$ prompts divide between benign requests phrased so as to look alarming and genuinely unsafe ones.
A response counts as a refusal when one of a fixed list of refusal phrases occurs in its first $120$ characters,
the same rule in every condition,
which keeps the measurement judge-free.
The two suites together separate a model that has become discriminating from one that declines whatever it is asked.
Refusal is reported alongside $Y$ and stays outside it (\apx{refusal}).

\paragraph{Entity leakage (memorization).}
Leakage is reported alongside $Y$ and stays outside it.
For the $375$ records that \textit{were} in the training set,
the model is shown the first half of the document and asked to continue it verbatim;
we count how many gold entities in the withheld second half are reproduced.
This is an extraction attack in miniature~\citep{carlini2021extracting} and serves as a manipulation check:
a mechanism whose leakage matches the baseline has not delivered the privacy whose alignment cost the paper measures.

\subsection{Semantic distortion}
The primary measure,
fixed before the main pass,
is the embedding distance of \S\ref{sec:design},
computed with \texttt{BAAI/\allowbreak bge-base-\allowbreak en-v1.5}~\citep{xiao2024cpack} as the sentence encoder.
$S$ is additionally reported as two components:
$S_{\mathrm{ent}}$,
the fraction of named entities and numbers not preserved between original and perturbed text,
defined against gold spans,
so reported on the entity-rich dataset;
and $S_{\mathrm{gen}}$,
the residual embedding drift after masking entities.
We test whether $S_{\mathrm{ent}}$ drives truthfulness and hallucination degradation while $S_{\mathrm{gen}}$ drives instruction-following degradation,
by refitting \eqref{eq:model} with the two components in place of $S$. Because the embedding-based components rest on a single encoder,
every distortion score is recomputed with a second,
unrelated encoder (\texttt{all-\allowbreak MiniLM-\allowbreak L6-v2}~\citep{wang2020minilm, reimers2019sbert}) as a robustness check;
the primary measure was fixed in advance so that it could not be selected after the results were visible.
The entity component is a token-membership rate rather than an embedding distance,
so it is invariant to that choice by construction.

\subsection{Experiment coverage}
Random replacement has four rates and private optimization four budgets;
sanitization has eight,
whose cost in collinearity exceeded its gain in coverage (\apx{deviations}).
Together the mechanisms give $17$ conditions per model on the entity-poor dataset:
the non-private baseline,
four rates of random replacement,
eight budgets of sanitization,
and four budgets of private optimization.
The entity-rich dataset adds two budgets of the entity-targeted control,
for $19$.
Across three base models and three seeds this is $153$ and $171$ adapters respectively,
$324$ in total,
of which $288$ enter \eqref{eq:model}:
the sixteen registered mechanism conditions per dataset.
The baseline sets the zero point,
and the entity-targeted control is a post-hoc addition analyzed separately (\S\ref{sec:carrier}).
Counting calibration,
the sensitivity blocks of \apx{sensitivity},
and repeated runs,
$444$ training runs were executed.
The three seeds supply the within-condition variance,
so every reported point carries an error bar.
